% ================================================
% DOCUMENT PRINCIPAL - COURS POO JAVA
% ================================================

% Inclure le préambule
% ================================================
% PRÉAMBULE COMPLET - COURS POO JAVA
% ================================================

\documentclass[12pt,a4paper]{article}

% ================================================
% PACKAGES ESSENTIELS
% ================================================

% Langue et encodage
\usepackage[utf8]{inputenc}
\usepackage[T1]{fontenc}
\usepackage[french]{babel}

% Mise en page
\usepackage[margin=2.5cm]{geometry}
\usepackage{fancyhdr}
\usepackage{lastpage}

% Mathématiques
\usepackage{amsmath, amssymb, amsthm}

% Graphiques et couleurs
\usepackage{graphicx}
\usepackage{xcolor}
\usepackage{tikz}

% Tableaux
\usepackage{array}
\usepackage{booktabs}
\usepackage{multirow}

% Listes
\usepackage{enumitem}

% Code source
\usepackage{listings}
\usepackage{minted} % Alternative (nécessite Python + Pygments)

% Boîtes colorées
\usepackage{tcolorbox}
\tcbuselibrary{listings, skins, breakable}

% Liens hypertextes
\usepackage{hyperref}

% Icônes
\usepackage{fontawesome5}

% QR Codes (optionnel)
\usepackage{qrcode}

% ================================================
% CONFIGURATION HYPERREF
% ================================================

\hypersetup{
    colorlinks=true,
    linkcolor=blue,
    filecolor=magenta,      
    urlcolor=cyan,
    pdftitle={Cours POO Java - Programmation Orientée Objet},
    pdfauthor={Votre Nom},
    pdfsubject={Programmation Orientée Objet en Java},
    pdfkeywords={Java, POO, Cours, Programmation},
    pdfcreator={LaTeX}
}

% ================================================
% COULEURS PERSONNALISÉES
% ================================================

\definecolor{githubblue}{RGB}{0, 119, 200}
\definecolor{githubgray}{RGB}{36, 41, 46}
\definecolor{javagreen}{RGB}{0, 128, 0}
\definecolor{javapurple}{RGB}{128, 0, 128}
\definecolor{javablue}{RGB}{0, 0, 255}
\definecolor{javared}{RGB}{255, 0, 0}
\definecolor{codebackground}{RGB}{248, 248, 248}

% ================================================
% CONFIGURATION LISTINGS (CODE JAVA)
% ================================================

\lstdefinestyle{javastyle}{
    language=Java,
    backgroundcolor=\color{codebackground},
    basicstyle=\ttfamily\small,
    keywordstyle=\color{javablue}\bfseries,
    stringstyle=\color{javared},
    commentstyle=\color{javagreen}\itshape,
    numberstyle=\tiny\color{gray},
    numbers=left,
    numbersep=8pt,
    tabsize=4,
    showspaces=false,
    showstringspaces=false,
    showtabs=false,
    breaklines=true,
    breakatwhitespace=true,
    frame=single,
    rulecolor=\color{gray},
    captionpos=b,
    xleftmargin=15pt,
    xrightmargin=5pt,
    aboveskip=10pt,
    belowskip=10pt
}

\lstset{style=javastyle}

% ================================================
% BOÎTES TCOLORBOX PERSONNALISÉES
% ================================================

% Boîte GitHub
\newtcolorbox{githubbox}[1]{
    colback=githubblue!5!white,
    colframe=githubblue,
    fonttitle=\bfseries,
    title={\faGithub\ #1},
    boxrule=0.8pt,
    arc=3mm,
    left=8pt,
    right=8pt,
    top=8pt,
    bottom=8pt,
    breakable
}

% Boîte Définition
\newtcolorbox{definition}[1]{
    colback=green!5!white,
    colframe=green!75!black,
    fonttitle=\bfseries,
    title={\faBook\ Définition : #1},
    boxrule=0.8pt,
    arc=3mm,
    breakable
}

% Boîte Remarque
\newtcolorbox{remarque}{
    colback=orange!5!white,
    colframe=orange!75!black,
    fonttitle=\bfseries,
    title={\faExclamationTriangle\ Remarque},
    boxrule=0.8pt,
    arc=3mm,
    breakable
}

% Boîte Exemple
\newtcolorbox{exemple}[1]{
    colback=blue!5!white,
    colframe=blue!75!black,
    fonttitle=\bfseries,
    title={\faLightbulb\ Exemple : #1},
    boxrule=0.8pt,
    arc=3mm,
    breakable
}

% Boîte Exercice
\newtcolorbox{exercice}[1]{
    colback=purple!5!white,
    colframe=purple!75!black,
    fonttitle=\bfseries,
    title={\faPencilAlt\ Exercice #1},
    boxrule=0.8pt,
    arc=3mm,
    breakable
}

% ================================================
% COMMANDES PERSONNALISÉES GITHUB
% ================================================

% URL du dépôt
\newcommand{\repourl}{https://github.com/GUEDDES/Cours-POO-Java}

% Lien GitHub simple
\newcommand{\githublink}[1]{%
    \href{\repourl/#1}{\texttt{#1}}%
}

% Lien GitHub avec icône
\newcommand{\githubrepo}{%
    \faGithub\ \href{\repourl}{github.com/GUEDDES/Cours-POO-Java}%
}

% Commande pour encart "Voir GitHub"
\newcommand{\voirgithub}[2]{%
    \begin{githubbox}{Ressources GitHub}
        \textbf{#1}
        
        \vspace{3pt}
        
        Voir sur GitHub : \href{\repourl/#2}{\texttt{#2}}
    \end{githubbox}
}

% Commande pour lien vers exemple
\newcommand{\githubexemple}[2]{%
    \href{\repourl/blob/main/#1}{\faCode\ #2}%
}

% Commande pour lien vers exercice
\newcommand{\githubexercice}[2]{%
    \href{\repourl/blob/main/#1}{\faPencilAlt\ #2}%
}

% ================================================
% EN-TÊTE ET PIED DE PAGE
% ================================================

\pagestyle{fancy}
\fancyhf{} % Clear default
\fancyhead[L]{\textbf{Cours POO Java}}
\fancyhead[R]{\leftmark}
\fancyfoot[L]{\githubrepo}
\fancyfoot[C]{Page \thepage\ sur \pageref{LastPage}}
\fancyfoot[R]{\today}
\renewcommand{\headrulewidth}{0.4pt}
\renewcommand{\footrulewidth}{0.4pt}

% ================================================
% COMMANDES UTILES
% ================================================

% Commande Java inline
\newcommand{\java}[1]{\texttt{#1}}

% Classe Java
\newcommand{\javaclass}[1]{\texttt{\textbf{#1}}}

% Méthode Java
\newcommand{\javamethod}[1]{\texttt{#1()}}

% ================================================
% TITRE DU DOCUMENT
% ================================================

\title{%
    \huge\textbf{Programmation Orientée Objet} \\[0.5cm]
    \Large Cours Complet avec Java \\[0.3cm]
    \large Année Universitaire 2024-2025
}

\author{Votre Nom}
\date{\today}

% ================================================
% FIN DU PRÉAMBULE
% ================================================


% ================================================
% DÉBUT DU DOCUMENT
% ================================================

\begin{document}

% ================================================
% PAGE DE TITRE
% ================================================

\maketitle
\thispagestyle{empty}

% ================================================
% ENCART GITHUB PRINCIPAL
% ================================================

\vspace{1cm}

\begin{tcolorbox}[colback=githubblue!5!white, colframe=githubblue, boxrule=1.5pt, arc=5mm]
    \centering
    \Large \faGithub\ \textbf{Dépôt GitHub du Cours}
    
    \vspace{8pt}
    
    \large Tous les exemples, exercices, diagrammes et solutions sont disponibles sur :
    
    \vspace{5pt}
    
    {\LARGE \href{\repourl}{\texttt{github.com/GUEDDES/Cours-POO-Java}}}
    
    \vspace{8pt}
    
    \begin{minipage}{0.6\textwidth}
        \textbf{Contenu du dépôt :}
        \begin{itemize}
            \item \faCode\ 9 exemples de code complets
            \item \faPencilAlt\ 8 exercices avec solutions
            \item \faProjectDiagram\ 3 diagrammes UML
            \item \faTerminal\ Sorties d'exécution
            \item \faBook\ Documentation complète
        \end{itemize}
    \end{minipage}
    \hfill
    \begin{minipage}{0.35\textwidth}
        \centering
        \qrcode[height=4cm]{\repourl}
        
        \small Scannez pour accéder
    \end{minipage}
\end{tcolorbox}

\vspace{1cm}

\begin{center}
    \textit{Ce cours est mis \`a jour régulièrement sur GitHub}
\end{center}

\newpage

% ================================================
% TABLE DES MATIÈRES
% ================================================

\tableofcontents
\newpage

% ================================================
% SECTION 1 : INTRODUCTION
% ================================================

\section{Introduction à la POO}
\label{sec:introduction}

\subsection{Qu'est-ce que la POO ?}

La Programmation Orientée Objet (POO) est un paradigme de programmation...

\voirgithub{Exercices Section 1}{01-Introduction/exercices}

% ================================================
% SECTION 2 : GÉNIE LOGICIEL
% ================================================

\section{Génie Logiciel}
\label{sec:genie-logiciel}

\subsection{Cycle de vie du logiciel}

Le développement d'un logiciel suit plusieurs phases...

\voirgithub{Exercices avec solutions}{02-Genie-Logiciel/exercices}

% ================================================
% SECTION 3 : BASES DE JAVA
% ================================================

\section{Bases de Java}
\label{sec:bases-java}

\subsection{Syntaxe de base}

Java utilise une syntaxe similaire au C++...

\begin{githubbox}{Exemples de Code}
    \textbf{\faCode\ Trois exemples disponibles}
    \begin{itemize}
        \item \githubexemple{03-Java-Bases/exemples/HelloWorld.java}{HelloWorld.java}
        \item \githubexemple{03-Java-Bases/exemples/CalculatriceSimple.java}{CalculatriceSimple.java}
        \item \githubexemple{03-Java-Bases/exemples/GestionTableaux.java}{GestionTableaux.java}
    \end{itemize}
\end{githubbox}

% ================================================
% SECTION 4 : CLASSES ET OBJETS
% ================================================

% Inclure la section complète depuis le fichier séparé
% ================================================
% EXEMPLE DE SECTION COMPLÈTE AVEC INTÉGRATION GITHUB
% Section 4 : Classes et Objets
% ================================================

\section{Classes et Objets}
\label{sec:classes-objets}

% ================================================
% INTRODUCTION
% ================================================

\subsection{Introduction}

La \textbf{Programmation Orientée Objet} (POO) repose sur le concept fondamental de \emph{classe} et d'\emph{objet}. Une classe est un modèle (ou \textit{blueprint}) qui définit les caractéristiques et comportements d'un type d'objet.

\begin{definition}{Classe}
    Une \textbf{classe} est une structure qui encapsule des données (attributs) et des comportements (méthodes) pour modéliser un concept du monde réel.
\end{definition}

\begin{definition}{Objet}
    Un \textbf{objet} est une instance concrète d'une classe. C'est une entité qui possède un état (valeurs des attributs) et un comportement (méthodes).
\end{definition}

% ================================================
% SYNTAXE
% ================================================

\subsection{Syntaxe d'une Classe}

En Java, une classe se définit avec le mot-clé \java{class} :

\begin{lstlisting}[language=Java, caption=Définition d'une classe Voiture]
public class Voiture {
    // Attributs (variables d'instance)
    private String marque;
    private String modele;
    private int annee;
    private double kilometrage;
    private boolean estDemarree;
    
    // Constructeur par défaut
    public Voiture() {
        this.marque = "Inconnue";
        this.modele = "Inconnu";
        this.annee = 2024;
        this.kilometrage = 0.0;
        this.estDemarree = false;
    }
    
    // Constructeur paramétré
    public Voiture(String marque, String modele, int annee) {
        this.marque = marque;
        this.modele = modele;
        this.annee = annee;
        this.kilometrage = 0.0;
        this.estDemarree = false;
    }
    
    // Méthodes (comportements)
    public void demarrer() {
        if (!estDemarree) {
            estDemarree = true;
            System.out.println(marque + " " + modele + " a démarré");
        }
    }
    
    public void arreter() {
        if (estDemarree) {
            estDemarree = false;
            System.out.println(marque + " " + modele + " s'est arrêtée");
        }
    }
    
    public void rouler(double distance) {
        if (estDemarree) {
            kilometrage += distance;
            System.out.printf("Vous avez parcouru %.1f km%n", distance);
        }
    }
}
\end{lstlisting}

% ================================================
% ENCART GITHUB
% ================================================

\vspace{10pt}
\begin{githubbox}{Ressources GitHub - Section 4}
    \textbf{\faCode\ Exemples de Code}
    \begin{itemize}
        \item \githubexemple{04-Classes-Objets/exemples/Voiture.java}{Voiture.java} - Classe complète avec tests
        \item \githubexemple{04-Classes-Objets/exemples/CompteBancaire.java}{CompteBancaire.java} - Gestion de compte bancaire
        \item \href{\repourl/blob/main/04-Classes-Objets/exemples/screenshots/output-examples.txt}{\faTerminal\ Sorties d'exécution}
    \end{itemize}
    
    \vspace{5pt}
    
    \textbf{\faPencilAlt\ Exercices Pratiques}
    \begin{itemize}
        \item \githubexercice{04-Classes-Objets/exercices/exercice1.md}{Énoncé : Classe Étudiant}
        \item \githubexemple{04-Classes-Objets/solutions/Exercice1Solution.java}{Solution complète}
    \end{itemize}
    
    \vspace{5pt}
    
    \textit{Accès direct :} \href{\repourl/tree/main/04-Classes-Objets}{\texttt{04-Classes-Objets/}}
\end{githubbox}

% ================================================
% UTILISATION
% ================================================

\subsection{Création et Utilisation d'Objets}

Pour utiliser une classe, on crée des \textbf{instances} (objets) :

\begin{lstlisting}[language=Java, caption=Utilisation de la classe Voiture]
public class TestVoiture {
    public static void main(String[] args) {
        // Création d'objets (instanciation)
        Voiture voiture1 = new Voiture();
        Voiture voiture2 = new Voiture("Peugeot", "308", 2023);
        
        // Utilisation des méthodes
        voiture2.demarrer();
        voiture2.rouler(50.5);
        voiture2.arreter();
    }
}
\end{lstlisting}

\begin{remarque}
    L'opérateur \java{new} alloue de la mémoire pour l'objet et appelle le constructeur correspondant.
\end{remarque}

% ================================================
% ENCAPSULATION
% ================================================

\subsection{Encapsulation}

L'\textbf{encapsulation} est un principe fondamental de la POO qui consiste \`a :
\begin{itemize}
    \item Rendre les attributs \java{private}
    \item Fournir des méthodes \java{public} pour y accéder (\textit{getters} et \textit{setters})
\end{itemize}

\begin{exemple}{Getters et Setters}
\begin{lstlisting}[language=Java]
public class Voiture {
    private String marque;
    
    // Getter
    public String getMarque() {
        return marque;
    }
    
    // Setter avec validation
    public void setMarque(String marque) {
        if (marque != null && !marque.isEmpty()) {
            this.marque = marque;
        } else {
            throw new IllegalArgumentException("Marque invalide");
        }
    }
}
\end{lstlisting}
\end{exemple}

% ================================================
% EXERCICES
% ================================================

\subsection{Exercices}

\begin{exercice}{1 - Classe Étudiant}
    Créer une classe \javaclass{Etudiant} avec les attributs suivants :
    \begin{itemize}
        \item \java{matricule} (String)
        \item \java{nom} et \java{prenom} (String)
        \item \java{moyenne} (double)
        \item \java{nbAbsences} (int)
    \end{itemize}
    
    Implémenter :
    \begin{itemize}
        \item Constructeurs (par défaut et paramétré)
        \item Getters et setters avec validation
        \item Méthode \javamethod{estAdmis} (retourne \java{true} si moyenne $\geq$ 10)
        \item Méthode \javamethod{afficher} (affiche toutes les informations)
    \end{itemize}
    
    \vspace{5pt}
    
    \textbf{\faGithub\ Solution :} \githubexemple{04-Classes-Objets/solutions/Exercice1Solution.java}{Voir sur GitHub}
\end{exercice}

% ================================================
% RÉSUMÉ
% ================================================

\subsection{Points Clés}

\begin{itemize}
    \item Une \textbf{classe} est un modèle, un \textbf{objet} est une instance
    \item Les \textbf{attributs} représentent l'état, les \textbf{méthodes} le comportement
    \item L'\textbf{encapsulation} protège les données
    \item Les \textbf{constructeurs} initialisent les objets
    \item Le mot-clé \java{this} référence l'objet courant
\end{itemize}

% ================================================
% FIN DE LA SECTION
% ================================================


% ================================================
% SECTIONS SUIVANTES (5-12)
% ================================================

\section{\u00c9léments Statiques}
\label{sec:elements-statiques}

% Contenu...

\voirgithub{Exemple : Classe Calculatrice}{05-Elements-Statiques/exemples/Calculatrice.java}

%------------------------------------------------

\section{Relations entre Classes}
\label{sec:relations}

% Contenu...

\begin{githubbox}{Exemples et Diagrammes}
    \textbf{\faProjectDiagram\ Modélisation}
    \begin{itemize}
        \item \href{\repourl/tree/main/06-Relations-Classes/diagrammes}{Diagrammes UML}
        \item \githubexemple{06-Relations-Classes/exemples/VoitureMoteur.java}{Composition : VoitureMoteur}
        \item \githubexemple{06-Relations-Classes/exemples/UniversiteEtudiant.java}{Agrégation : Université}
    \end{itemize}
\end{githubbox}

%------------------------------------------------

\section{Héritage}
\label{sec:heritage}

% Contenu...

\voirgithub{Exemple : Hiérarchie Animal}{09-Heritage/exemples/AnimalHierarchie.java}

%------------------------------------------------

\section{Polymorphisme}
\label{sec:polymorphisme}

% Contenu...

\voirgithub{Exemple : Système de Paiement}{10-Polymorphisme/exemples/SystemePaiement.java}

%------------------------------------------------

\section{Exceptions}
\label{sec:exceptions}

% Contenu...

\voirgithub{Exemple : Gestion de Fichiers}{12-Exceptions/exemples/GestionFichiers.java}

% ================================================
% CONCLUSION
% ================================================

\section{Conclusion}

Ce cours a couvert les concepts fondamentaux de la POO en Java...

\vspace{10pt}

\begin{center}
    \Large \textbf{Pour aller plus loin}
    
    \vspace{5pt}
    
    \githubrepo
\end{center}

% ================================================
% FIN DU DOCUMENT
% ================================================

\end{document}
